\section{Introduction}
Information security risk management is a central part of modern management and control systems in businesses. Standards such as ISO/IEC 27005 and frameworks such as NIST SP 800-37 describe structured processes for identifying, analyzing, and treating risk.

At the same time, businesses have become increasingly dependent on digital systems, increasing both complexity and vulnerability. This has made systematic risk management a critical part of both operational operations and strategic management.

In this context, it is relevant to examine the extent to which existing risk management tools support requirements for traceability, documentation, analysis, and decision-making, as well as what strengths and limitations can be identified when using such tools.

To understand how risk management can be operationalized in practice, it is appropriate to distinguish between methodology and tool support. The next section explains what risk management entails and what types of tools can be used to support the process
\subsection{Project background}
% TODO
Risk management involves being able to; identify threats and vulnerabilities, assess probability and consequence, and implement measures to reduce risk. Methodology describes how the process should be carried out, tools support the implementation itself, and contribute to traceability and consistency. 

There are various tools that can be used to achieve this:
\begin{itemize}
    \item Simple tools that can be used: excel or some other similar spreadsheet program. These provide basic support, but are more limited in traceability

    \item Open source tools: For example MONARC, Risk BowTie (Legg til de vi har) offer structured methodological support at no cost.

    \item Commercial GRC solutions: (Legg til de vi har fått tak i) provide many functionalities, can be seamlessly integrated into the system, and provide advanced reporting.
\end{itemize}
The choice of these available solutions will be influenced by the size of the business, the need for documentation and analysis. The next steps are to assess whether these available tools actually meet the requirements of ISO 27005, NIST RMF, and the ENISA recommendations.

\subsection{Project description}
% TODO
This project aims to evaluate different tools used for risk management within information security. We will combine theoretical analysis with established standards and frameworks based on selected tools. 

The study will be organized into five phases. 
\begin{itemize}
    \item Phase one will be a literature review based on ENISA risk management reports, ISO/IEC 27005, NIST SP 800-73 (RMF), OCTAVE, FAIR, and other relevant publications on GRC and cyber risk quantification. This phase is to identify requirements and establish a structured set of evaluations criteria.

    \item Phase two - an overview of 8-15 relevant risk management tools will be developed. The tools will consist of open source, commercial platforms, and semi-comercial or academic tools.
    
    \item Phase three - a structured evaluation framework will be defined. The criteria will include alignment with established risk management frameworks, functional capabilities, integration options, usability, and administrative aspects.

    \item Phase four - Select four to six tools that will be used in an in-depth analysis. (hva velger vi å utføre her?)

    \item Phase five - The results gained from our analysis will be used to do a comperative assesment to identify their strengths, weaknesses, and practical suitability across small and medium-sized enterprises, large enterprises, and public sector organizations. 
\end{itemize}
 

\subsection{Project scope}
% TODO

\subsection{Research question}
% TODO

\subsection{Project goals}
% TODO

\subsection{Thesis structure}
% TODO